%% ma: L12-B08-0018
%% lop: 12
%% chuong: 2
%% bai: 8
%% dang: 12.8.5
%% loai: TLN
%% muc: VD
%% dap_an: "42,5"
%% nguon: "(NBV)-ĐỀ SỐ 10-MỖI NGÀY 1 ĐỀ THI 2025.pdf (D10, MỖI NGÀY 1 ĐỀ - Đề số 10), Phần III Câu 3"
%% kien_thuc: "Bài 8 (vectơ cùng phương, tọa độ), Bài 7 (hệ trục tọa độ, mặt phẳng Oxy)"
%% trang_thai: chinh-thuc
%% ghi_chu: "Cùng ý tưởng với L12-B08-0005 (đường thẳng qua hai điểm cắt mặt phẳng Oxy, dùng điều kiện thẳng hàng), nhập vào dạng 12.8.5 đã duyệt. Hình gốc là ảnh máy bay trang trí nên không vẽ hình."
\begin{ex}
Khi gắn hệ tọa độ $Oxyz$ (đơn vị trên mỗi trục tính theo kilômét) vào một sân bay, mặt phẳng $(Oxy)$ trùng với mặt sân bay. Một máy bay bay theo đường thẳng từ vị trí $A(5;0;5)$ đến vị trí $B(10;10;3)$ và hạ cánh tại vị trí $M(a;b;0)$. Giá trị của $a+b$ bằng bao nhiêu (viết kết quả dưới dạng số thập phân)?
\shortans{42,5}
\loigiai{Ta có $\overrightarrow{AB}=(5;10;-2)$ và $\overrightarrow{AM}=(a-5;b;-5)$. Ba điểm $A$, $B$, $M$ thẳng hàng nên $\overrightarrow{AM}=k\overrightarrow{AB}$, tức $(a-5;b;-5)=k(5;10;-2)$. Từ $-5=-2k$ suy ra $k=\dfrac52$, do đó $a=5+\dfrac{25}2=17{,}5$ và $b=25$. Vậy $a+b=42{,}5$.}
\end{ex}
